\section*{Chapitre \ref{ch:hf:inventory-models} --- Modèles d'inventaire}

\begin{solution}{exo:hf:inventory-models:1}
$S^{\mathrm{res}}=100-3\times0.1\times4\times0.5=99.4$.
\end{solution}

\begin{solution}{exo:hf:inventory-models:2}
$\Delta=0.1\times4\times1+20\ln(1+0.1/1.5)=0.40+1.29=1.69$~; la partie monopolistique vaut 1.29.
\end{solution}

\begin{solution}{exo:hf:inventory-models:3}
$\ln(40/18)=0.5k$ donne $k=1.60$~; $A=0.04\,e^{0.5k}=0.0889$ exécution par seconde.
\end{solution}

\begin{solution}{exo:hf:inventory-models:4}
À l'horizon, le terme d'inventaire s'annule et $(2/\gamma)\ln(1+\gamma/k)\to2/k$ lorsque $\gamma\to0$ (puisque $\ln(1+x)\sim x$). Un teneur de marché neutre au risque
arbitre encore entre le gain par exécution, qui croît avec la profondeur, et le taux d'exécution, qui décroît avec elle~: $\delta Ae^{-k\delta}$ est maximisé en
$\delta=1/k$.
\end{solution}

\begin{solution}{exo:hf:inventory-models:5}
Entre $q=3$ (profondeur au cours vendeur 0.011) et $q=4$ ($-0.217$). Une profondeur négative offre sous le mid~: dans un carnet réel elle croise la fourchette, c'est-à-dire que le teneur de
marché vend agressivement au cours acheteur.
\end{solution}

\begin{solution}{exo:hf:inventory-models:6}
La stratégie d'inventaire pratique un décalage des cotations~: dès qu'elle détient une position, l'une de ses cotations est plus proche du mid que la moitié de la largeur moyenne, et c'est le
côté qui s'exécute, donc il gagne moins par exécution (il s'exécute un peu plus souvent, 97.0 fois contre 91.7, pour moins d'argent). Ce qu'il gagne, c'est une variance réduite~:
6.49 contre 14.00. Pour toute aversion au risque positive, son équivalent certain est le plus élevé, et c'est celui que la stratégie maximise.
\end{solution}

\begin{solution}{exo:hf:inventory-models:7}
$q=0$~: les deux profondeurs valent 0.804. $q=3$~: cours acheteur 1.550, cours vendeur 0.011.
\end{solution}

\begin{solution}{exo:hf:inventory-models:8}
Les cotations placées derrière le touch d'une action cotée à un tick ne s'exécutent presque jamais (le Quoter explorateur du chapitre n'en a obtenu aucune en 5\,280 secondes)~: l'ajustement porte sur des zéros,
et $k$ n'est pas identifié. La profondeur ne peut pas être choisie continûment dans un carnet à grand tick~; la décision consiste à savoir s'il faut se placer au touch.
\end{solution}

\begin{solution}{pb:hf:inventory-models:1}
\begin{enumerate}
\item $S_t=S_0+\sigma W_t$~; cotations $S\mp\delta^{b,a}$~; \omterm{def:hf:inventory-models:intensity}{intensité d'exécution} $Ae^{-k\delta}$~; maximiser $\E[-e^{-\gamma(X_T+q_TS_T)}]$.
\item Le taux auquel une cotation au repos à la profondeur $\delta$ est exécutée~; $\lambda(\delta)=Ae^{-k\delta}$.
\item Le prix d'indifférence d'une unité de plus ou de moins~; $S^{\mathrm{res}}=S-q\gamma\sigma^2(T-t)$.
\item $\gamma\sigma^2(T-t)$ (risque d'inventaire) plus $(2/\gamma)\ln(1+\gamma/k)$ (partie monopolistique).
\item $S^{\mathrm{res}}=99.40$~; largeur $0.2+1.291=1.491$~; cours acheteur 98.65, cours vendeur 100.15.
\item Un coût $\phi\int q^2dt+\alpha q_T^2$ sur l'inventaire détenu~; $\phi=\tfrac12\gamma\sigma^2$ pénalise la variance que la position ajoute.
\item Voir la démonstration de la \cref{prop:hf:inventory-models:cj}.
\item La valeur est linéaire en $s$ ($x+qs+h$), donc la dérivée seconde en $s$ est nulle~; le risque n'intervient que par la pénalité.
\item 0.804 des deux côtés pour $q=0$~; cours acheteur 1.550 et cours vendeur 0.011 pour $q=3$~; au-delà de $q=3$, le cours vendeur est sous le mid, c'est-à-dire un ordre croisé.
\item $S_0=100$, $T=1$, $\sigma=2$, pas 0.005, $k=1.5$, $A=140$, mid binomial, 1\,000 trajectoires.
\item $\gamma=0.01$~: 68.45 (9.06) contre 68.50 (13.91)~; $\gamma=0.1$~: 64.83 (6.49) contre 68.15 (14.00)~; $\gamma=1$~: 31.44 (4.83) contre 43.30
(10.34)~; l'article~: 68.6 (8.7)/68.8 (12.8), 65.0 (6.6)/68.4 (12.7), 31.4 (5.0)/44.0 (11.0).
\item À partir de $\gamma\approx0.1$~: 5~\% de gain en moins à 0.1, 27~\% en moins à 1~; à 0.01 les deux gains sont égaux.
\item Les cotations pénalisées exactes dominent les largeurs symétriques~: $\phi=0.5$ donne 66.7 avec un écart type de 6.9 contre 68.5 avec un écart type de 13.9.
\item Huit exécutions en 1\,640 secondes au touch et aucune en 5\,280 secondes derrière lui~: ici, les ordres au marché ne vont jamais au-delà du touch.
\item $\partial_th-\phi q^2+\lambda\max(0,c+h(q-1)-h(q))+\lambda\max(0,c+h(q+1)-h(q))=0$~; laisser une cotation d'un côté lorsque son crochet est positif.
\item À partir des 241 exécutions du Quoter symétrique sur trois sessions de calibration~: $\lambda=0.033$ lots par seconde et par côté, $c=0.139$ ticks par action~; des sessions
séparées, afin que la comparaison ne soit pas évaluée sur les données qui ont fixé ses paramètres.
\item Sur le marché simulé~: symétrique $-\$65.17$ avec une moyenne de $|q|$ de 2.51 lots~; $\phi=10^{-3}$~: $+\$11.00$ avec 0.57. Dans le modèle à $\gamma=0.1$~:
64.83 avec un écart type de 6.49 contre 68.15 avec un écart type de 14.00.
\item Les erreurs types (environ 40~\$ pour le Quoter symétrique) sont supérieures à certaines des différences, mais chaque session pénalisée détient moins
d'inventaire et varie moins.
\item L'absence d'information~: les exécutions sont suivies de mouvements défavorables, que le mid, martingale dans le modèle, exclut.
\item Cesser d'alourdir la position et continuer à afficher l'autre côté, à proportion du coût de détention de la position.
\end{enumerate}
\end{solution}

\begin{solution}{iq:hf:inventory-models:1}
Les deux, à la baisse~: un cours vendeur plus bas est plus souvent exécuté et un cours acheteur plus bas l'est moins. Sur un marché à un tick, il faut garder le cours vendeur au touch et retirer ou réduire le cours acheteur.
\iqlookfor{les deux cotations suivent le \omterm{def:hf:inventory-models:reservation}{prix de réservation}~; la version à grand tick.}
\end{solution}

\begin{solution}{iq:hf:inventory-models:2}
Maximiser $\delta\,Ae^{-k\delta}$~: $\delta^\ast=1/k$, soit un gain de $A/(ek)$ par unité de temps.
\iqlookfor{la condition du premier ordre et la valeur.}
\end{solution}

\begin{solution}{iq:hf:inventory-models:3}
Le risque de détenir $q$ jusqu'à l'horizon est $q^2\sigma^2(T-t)$~; il diminue à mesure que l'horizon approche. Un teneur de marché qui conserve ses positions pendant la nuit n'a pas de
horizon naturel $T$~: il utilise une pénalité continue (une charge de risque constante par unité de temps) ou un horizon égal à sa durée de détention habituelle.
\iqlookfor{l'origine du $T-t$ et un remplacement raisonnable.}
\end{solution}

\begin{solution}{iq:hf:inventory-models:4}
Sous-dépassement~: $e^{-\alpha kq^2}$ et les petites entrées de l'exponentielle de matrice deviennent nulles, ce qui donne $\ln0=-\infty$. Renormaliser $\omega$ à chaque pas (un
facteur constant disparaît de toutes les profondeurs) et borner le logarithme par le bas~; ou travailler dans l'espace logarithmique.
\iqlookfor{reconnaître un sous-dépassement et pourquoi le changement d'échelle est sans conséquence.}
\end{solution}

\begin{solution}{iq:hf:inventory-models:5}
Choisir $\phi$ de sorte que l'inventaire accepté par la politique, multiplié par la volatilité sur la période de détention, respecte la limite de perte à un niveau de confiance donné~; vérifier
cela en simulant la distribution du P\&L de la politique, et plafonner aussi l'inventaire de manière stricte.
\iqlookfor{relier la pénalité à un budget de risque et vérifier par simulation.}
\end{solution}

\begin{solution}{iq:hf:inventory-models:6}
Le choix des côtés à afficher au touch (et de la taille, et du moment où croiser la fourchette). L'état est $(t,q)$ et peut-être la position dans la file~; le contrôle de chaque
côté est binaire, l'\omterm{def:hf:inventory-models:intensity}{intensité d'exécution} est celle du touch, et l'équation de Hamilton--Jacobi--Bellman admet un maximum entre afficher la cotation et s'abstenir.
\iqlookfor{contrôles binaires et HJB à choix discret.}
\end{solution}
